\addcontentsline{toc}{section}{Введение}
\section*{Введение}

В данной лабораторной работе необходимо разработать функциональную схему электронных часов в соответствии с вариантом \textbf{101101132}. Базовые функции часов включают отображение и корректировку текущего времени.

Дополнительные функции часов, заданные вариантом:
\begin{itemize}
    \item A = 1 -- отображение и корректировка секунд;
    \item B = 0 -- дополнительные счётчики времени в часовых поясах отсутствуют;
    \item C = 1 -- раздельная корректировка десятков и единиц разрядов;
    \item D = 1 -- 24-часовой режим работы;
    \item E = 0 -- отключение индикаторов для экономии электроэнергии отсутствует;
    \item F = 1 -- останов часов по нажатию кнопки;
    \item G = 1 -- простой секундомер с возможностью сброса, запуска и остановки;
    \item H = 3 -- звуковой сигнал каждый час с выбором мелодии по номеру;
    \item I = 2 -- будильник на устанавливаемое время с выбором мелодии по номеру и возможностью отключения.
\end{itemize}

Этапы выполнения работы:

\begin{enumerate}
\item Построить граф управляющего автомата часов и дать пояснения к нему. Пояснения предполагают
описание логического смысла каждого состояния, перечень визуальной информации, выводимой на
индикаторы, а также порядок использования всех тех возможностей часов, которые перечислены в задании.
\item Изобразить общую структурную схему электронных часов с указанием всех необходимых управляющих
микрокоманд (импульсных и потенциальных). Функции каждого блока структурной схемы должны быть
пояснены. Должны быть даны также пояснения функции всех управляющих микрокоманд.
\item Провести кодирование входных и выходных воздействий и состояний автомата.
\item Построить минимизацию функций блоков F и Fi.
\item Построить общую функциональную схему. При этом необходимо четко описать алгоритм работы и уметь
объяснить принцип проектирования всех блоков.
\item Определить (приблизительно) площадь микросхемы, реализующей построенную функциональную схему
при современной плотности компоновки транзисторов
\end{enumerate}